%!TEX root = iopjournal.tex
%
% Evaluation and results
%
% Inhalt (Vorschlag):
%   Bewertungsprotokoll
%     - Simulationsfehler als Bewertungsgroesse, Normierung am nominalen Lauf
%     - Ground Truth der injizierten Drifts, Toleranz und Cooldown
%     - Score-Funktionen fuer Detektion und Adaption
%   Ergebnisse
%     - drift-bedingte Modellverschlechterung: RMSE, Streuung, Fehlerverhaeltnis,
%       Aufteilung nach Kampagne und uebrigen Betriebspunkten
%     - Abstand zum naechsten Kampagnenpunkt: das Fehlerniveau entscheidet ueber
%       die Sichtbarkeit des Drifts, nicht dessen Groesse
%     - Detektion: KSWIN, ADWIN, DDM, EDDM auf dem Simulationsfehler (frouros)
%     - Adaption: passiv, aktiv, kombiniert; inkrementelles Nachtraining
%
% Aus Abschnitt 3 offen:
%   - Schwelle zur Binarisierung des Regressionsfehlers fuer DDM/EDDM berichten;
%     in 3.2 als Parameter eingefuehrt
%   - Schwellwert auf der rollenden RMSE als Referenzdetektor mitfuehren; in 3.2
%     als einfachste Form angekuendigt
%   - Verzugszeit als Kennzahl an das Driftmuster aus 3.1 binden
%   - kadlec.2011 (S. 3) aufgreifen: keine Adaption, solange der Hardware-Sensor
%     falsch misst -- der eigene Fall unterscheidet sich, weil das Label sauber
%     bleibt und nur ein Feature verfaelscht ist
%   - der Drift wird nur entlang der gefahrenen Trajektorie sichtbar, wirkt aber
%     im ganzen Bereich -> Begruendung fuer den nominalen Referenzlauf
%
% Vgl. Dissertation: sec:tep_modellierung, sec:tep_detektion, sec:tep_adaption
%

\section{Evaluation and results} \label{sec:evaluation}


%
% ---------------------------------------------------------------------------
% Beispielcode aus dem IOP-Template, hier als Vorlage abgelegt.
% Hinweis: figure1.pdf liegt nicht mehr im Verzeichnis.
% ---------------------------------------------------------------------------
%
% \begin{figure}
%  \centering
%         \includegraphics[width=0.5\textwidth]{figure1}
%  \caption{Text describing the figure and the main conclusions drawn from it.}
% \label{fig1}
% \end{figure}
%
% \begin{table}
% \caption{Caption text describing the table.}
% \centering
% \begin{tabular}{l c c c}
% \hline
% Column heading & Column heading & Column heading & Column heading \\
% \hline
% Data row 1 & 1.0 & 1.5 & 2.0 \\
% \hline
% \end{tabular}
% \label{tab1}
% \end{table}
%
